\section{Polygamma functions and generalized harmonic sums}
\label{sec:harmonics}

\subsection{Uncompensated polygamma pairings}
For uncompensated Hurwitz zeta functions, the Laplace kernel $Q(x)$ has a
simple pole at zero. Consequently, $\Rea s,\Rea t>3/2$ is a sufficient
$L^2$ region for the same argument to give
\begin{equation}\label{eq:plain-zeta-pair}
 \frac1{2\pi}\int_{\R}\zeta(s,a+\ii y)\zeta(t,b-\ii y)\dd y
 =\frac{\Gamma(w)}{\Gamma(s)\Gamma(t)}
       \{\zeta(w-1,A)+(1-A)\zeta(w,A)\}.
\end{equation}
This follows directly from the $Q^2$ calculation in the proof of
\cref{thm:basic}; no compensation is needed in this smaller domain.
Using the standard identity
$\psi^{(r)}(z)=(-1)^{r+1}r!\zeta(r+1,z)$ gives the following consequence.

\begin{corollary}\label{thm:polygamma}
For integers $r,k\ge1$, $\Rea a,\Rea b>0$, and $A=a+b$,
\begin{align}
 &\frac1{2\pi}\int_{\R}\psi^{(r)}(a+\ii y)
                                  \psi^{(k)}(b-\ii y)\dd y\notag\\
 &\quad=(-1)^{r+k}(r+k)!
      \{\zeta(r+k,A)+(1-A)\zeta(r+k+1,A)\}.\label{eq:polygamma-pair}
\end{align}
In particular,
\begin{equation}\label{eq:polygamma-norm}
 \frac1{2\pi}\int_{\R}
       |\psi^{(r)}(\tfrac12+\ii y)|^2\dd y=(2r)!\zeta(2r).
\end{equation}
\end{corollary}
The restriction $r,k\ge1$ matters. Setting both to zero would introduce
the divergent squared norm of the uncorrected digamma function, not the
finite value in \eqref{eq:headline-digamma}.

\subsection{Finite generalized harmonic families}
For a positive integer $N$, define
\begin{equation}\label{eq:finite-H}
 \mathcal H_N^{(s)}(z)=\sum_{j=0}^{N-1}(z+j)^{-s}.
\end{equation}
At $z=1$ this is the usual generalized harmonic number $H_N^{(s)}$.
For $\Rea s>0$, it is the Laplace transform with kernel
$\sum_{j=0}^{N-1}e^{-jx}$, divided by $\Gamma(s)$.

\begin{theorem}[Finite harmonic and mixed pairings]\label{thm:harmonic}
Let $M,N\ge1$, $\Rea a,\Rea b>0$, $\Rea s,\Rea t>1/2$, and
$A=a+b$, $w=s+t-1$. Then
\begin{align}
 &\frac1{2\pi}\int_{\R}\mathcal H_M^{(s)}(a+\ii y)
                           \mathcal H_N^{(t)}(b-\ii y)\dd y\notag\\
 &\qquad=\frac{\Gamma(w)}{\Gamma(s)\Gamma(t)}
       \sum_{j=0}^{M-1}\sum_{k=0}^{N-1}(A+j+k)^{-w},\label{eq:HH-pair}\\
 &\frac1{2\pi}\int_{\R}Z_s(a+\ii y)
                           \mathcal H_N^{(t)}(b-\ii y)\dd y\notag\\
 &\qquad=\frac{\Gamma(w)}{\Gamma(s)\Gamma(t)}
                         \sum_{k=0}^{N-1}Z_w(A+k).\label{eq:ZH-pair}
\end{align}
The second expression is taken at its removable value if $w=1$.
Arbitrary spectral and position derivatives are permitted inside both
identities in the stated domain.
\end{theorem}
\begin{proof}
Multiply the two finite exponential kernels and apply
\cref{thm:plancherel}. Each exponential term contributes
$\Gamma(w)(A+j+k)^{-w}$. In the mixed product it contributes
$\Gamma(w)Z_w(A+k)$ by \eqref{eq:Z-laplace}.
The kernels are bounded at zero, so the same $L^2$ holomorphy argument
applies to all derivatives.
\end{proof}

For $a+b=1$ and $s=t=1$, these become particularly elementary:
\begin{align}
 &\frac1{2\pi}\int_{\R}\mathcal H_M^{(1)}(a+\ii y)
                        \mathcal H_N^{(1)}(b-\ii y)\dd y\notag\\
 &\quad=(M+N)H_{M+N}-MH_M-NH_N,\label{eq:HH-elementary}\\
 &\frac1{2\pi}\int_{\R}G_0(a+\ii y)
                        \mathcal H_N^{(1)}(b-\ii y)\dd y\notag\\
 &\quad=\log\Gamma(N+1)-NH_N+N(1+\gamma).
               \label{eq:GH-elementary}
\end{align}
Indeed, the first double sum can be summed along its diagonals, and the
second is $\sum_{k=1}^N(\log k-\psi(k))$.
Taking spectral derivatives supplies finite sums of
$\log^r(A+j+k)/(A+j+k)^w$ and corresponding exact mixed identities for
higher $G_m$, with the same Gamma-quotient coefficient bookkeeping as
in \cref{sec:stieltjes}.

\section{Lerch products and polylogarithm order derivatives}
\label{sec:lerch}

Let $\xi$ lie in $\{\xi:|\xi|\le1,\ \xi\ne1\}$. The Lerch transcendent
has the integral representation \cite{dlmf-lerch}
\begin{equation}\label{eq:lerch-laplace}
 \Phi(\xi,s,z)=\frac1{\Gamma(s)}\int_0^\infty
           \frac{x^{s-1}e^{-zx}}{1-\xi e^{-x}}\dd x,
 \qquad\Rea s>0,\quad\Rea z>0.
\end{equation}
For $|\xi|<1$ this also follows by summing its defining series.
The kernel is bounded at zero and along the positive real axis. Boundary
colors $|\xi|=1$, $\xi\ne1$, follow by dominated continuation from inside
the disk. No assertion here includes the singular color $\xi=1$ without
the separate compensation already developed.

\begin{theorem}[Colored pairings and confluent colors]\label{thm:lerch}
Let $\Rea s,\Rea t>1/2$, $\Rea a,\Rea b>0$, and set
$A=a+b$, $w=s+t-1$. For admissible colors $\xi\ne\eta$,
\begin{align}
 &\frac1{2\pi}\int_{\R}\Phi(\xi,s,a+\ii y)
                             \Phi(\eta,t,b-\ii y)\dd y\notag\\
 &\quad=\frac{\Gamma(w)}{\Gamma(s)\Gamma(t)}
           \frac{\xi\Phi(\xi,w,A)-\eta\Phi(\eta,w,A)}{\xi-\eta}.
               \label{eq:lerch-distinct}
\end{align}
At $\eta=\xi$, the quotient is replaced by
\begin{equation}\label{eq:lerch-confluent}
 \Phi(\xi,w-1,A)+(1-A)\Phi(\xi,w,A).
\end{equation}
In particular, when $A=1$ and $\xi\ne\eta$,
\begin{align}
 &\frac1{2\pi}\int_{\R}\Phi(\xi,s,a+\ii y)
                             \Phi(\eta,t,1-a-\ii y)\dd y\notag\\
 &\quad=\frac{\Gamma(w)}{\Gamma(s)\Gamma(t)}
                 \frac{\Li_w(\xi)-\Li_w(\eta)}{\xi-\eta}.
                     \label{eq:polylog-divdiff}
\end{align}
For equal nonzero colors the final quotient is $\Li_{w-1}(\xi)/\xi$;
at $\xi=0$ its removable value is $1$.
\end{theorem}
\begin{proof}
For distinct colors, use the exact partial fraction identity
\[
 \frac1{(1-\xi r)(1-\eta r)}
 =\frac{\xi/(1-\xi r)-\eta/(1-\eta r)}{\xi-\eta}
\]
in the Laplace product supplied by \cref{thm:plancherel}.
For equal colors, the kernel has series
$\sum_{n\ge0}(n+1)\xi^n e^{-nx}$, and
$n+1=(n+A)+(1-A)$ proves \eqref{eq:lerch-confluent}.
This series argument initially uses $|\xi|<1$; its identity extends to
the stated boundary colors by the same dominated kernel argument.
Finally, $\xi\Phi(\xi,w,1)=\Li_w(\xi)$, including the removable
interpretation at zero. All products converge by $L^2$.
\end{proof}

Order differentiation of \eqref{eq:polylog-divdiff} is a direct source of
identities for $\partial_w^r\Li_w(\xi)$ and mixed order derivatives of
Lerch functions. In particular, at $s=t=1$, the total-order derivatives
are generated by
\[
 C(u,v)\frac{\Li_{1+u+v}(\xi)-\Li_{1+u+v}(\eta)}{\xi-\eta}.
\]
This is a convergent parameter identity near the origin, rather than an
experimental relation inferred from numerical integer-relation searches.

\subsection{Mixed Stieltjes--Lerch closure and the spectral resonance}
\begin{theorem}\label{thm:mixed-lerch}
Under the hypotheses of \cref{thm:lerch},
\begin{align}
 &\frac1{2\pi}\int_{\R}Z_s(a+\ii y)\Phi(\xi,t,b-\ii y)\dd y\notag\\
 &\quad=\frac{\Gamma(w)}{\Gamma(s)\Gamma(t)}
  \left\{\frac{\zeta(w,A)-\xi\Phi(\xi,w,A)}{1-\xi}
                     -\frac{\Phi(\xi,w-1,A)}{w-1}\right\}.
                         \label{eq:mixed-lerch}
\end{align}
The brace is holomorphic for $\Rea w>0$. At $s=t=1$ it gives
\begin{align}
 &\frac1{2\pi}\int_{\R}G_0(a+\ii y)\Phi(\xi,1,b-\ii y)\dd y\notag\\
 &\quad=\frac{-\psi(A)-\xi\Phi(\xi,1,A)}{1-\xi}
                                  -\Phi'(\xi,0,A).
                                      \label{eq:mixed-lerch-resonance}
\end{align}
\end{theorem}
\begin{proof}
The product kernel is
\[
 \frac{h(x)}{1-\xi e^{-x}}
 =\frac{Q(x)-\xi/(1-\xi e^{-x})}{1-\xi}
                     -\frac1{x(1-\xi e^{-x})}.
\]
Integrate the separated expression first for $\Rea w>1$.
The combined kernel is bounded at zero, so its Mellin transform is
holomorphic for $\Rea w>0$. At $w=1$, both singular summands have residue
$1/(1-\xi)$ with opposite signs, because
$\Phi(\xi,0,A)=1/(1-\xi)$. Their constant terms give
\eqref{eq:mixed-lerch-resonance}. Derivatives of this removable identity
are justified by the original holomorphic $L^2$ pairing.
\end{proof}

When $A=1$, \eqref{eq:mixed-lerch-resonance} becomes
\begin{equation}\label{eq:mixed-Li0}
 \frac1{2\pi}\int_{\R}G_0(a+\ii y)\Phi(\xi,1,1-a-\ii y)\dd y
 =\frac{\gamma+\log(1-\xi)}{1-\xi}
       -\frac{\left.\partial_s\Li_s(\xi)\right|_{s=0}}{\xi}.
\end{equation}
At zero color the final quotient is interpreted by continuity. At
$\xi=-1$, the right-hand side reduces to
\begin{equation}\label{eq:alternating-special}
 \frac\gamma2+\log2-\frac12\log\pi.
\end{equation}
Indeed, $\Li_s(-1)=(2^{1-s}-1)\zeta(s)$, whence
$\Li'_0(-1)=\frac12\log(2/\pi)$. Moreover,
\[
 \Phi(-1,1,z)=\frac12\left\{\psi\left(\frac{z+1}{2}\right)
                           -\psi\left(\frac z2\right)\right\}.
\]
Thus \eqref{eq:alternating-special} is also an explicit compensated
digamma--digamma-difference integral.

\subsection{Every nontrivial root of unity gives a finite Gamma formula}
If $\xi^q=1$ and $\xi\ne1$, grouping the Lerch series modulo $q$ gives
\begin{equation}\label{eq:root-unity}
 \Phi(\xi,s,A)=q^{-s}\sum_{r=0}^{q-1}\xi^r
                                \zeta\left(s,\frac{A+r}{q}\right).
\end{equation}
The zeta residues at $s=1$ cancel because $\sum\xi^r=0$. With
$c_r=(A+r)/q$, its two required values are
\begin{align}
 \Phi(\xi,1,A)&=-\frac1q\sum_{r=0}^{q-1}\xi^r\psi(c_r),\label{eq:root-one}\\
 \Phi'(\xi,0,A)&=\sum_{r=0}^{q-1}\xi^r\log\Gamma(c_r)
                                      -\frac{\log q}{1-\xi}.
                                          \label{eq:root-zero-derivative}
\end{align}
Substituting into \eqref{eq:mixed-lerch-resonance} yields the finite family
\begin{align}
 &\frac1{2\pi}\int_{\R}G_0(a+\ii y)\Phi(\xi,1,b-\ii y)\dd y\notag\\
 &\quad=\frac{-\psi(A)+(\xi/q)\sum_{r=0}^{q-1}\xi^r\psi(c_r)
                      +\log q}{1-\xi}
             -\sum_{r=0}^{q-1}\xi^r\log\Gamma(c_r).
                 \label{eq:root-Gamma-pair}
\end{align}
For higher indices, differentiating the finite sum
\eqref{eq:root-unity} gives a triangular combination of Hurwitz order
jets and powers of $\log q$. This is a proved finite reduction at every
index, with no presumption of further transcendence reductions.
